\section{Phương pháp nghiên cứu}
Khóa luận sử dụng phương pháp nghiên cứu kết hợp giữa khảo sát tài liệu, thiết kế hệ thống và triển khai thực nghiệm. Trước hết, các thông số liên quan đến chất lượng không khí trong nhà và điều kiện sinh trưởng của cây được tổng hợp từ các tài liệu về IAQ, IoT và hệ thống trồng cây/thủy canh thông minh \cite{who_air_pollution,epa_iaq,bose_iot_hydroponics}. Trên cơ sở đó, đề tài xác định nhóm cảm biến cần triển khai gồm CO2, nhiệt độ, độ ẩm, ánh sáng, TDS và pH. Kiến trúc phần cứng và phần mềm sau đó được thiết kế theo hướng phân lớp, trong đó node đảm nhiệm thu thập dữ liệu, gateway thực hiện truyền thông lên nền tảng IoT, ThingsBoard quản lý dữ liệu và ứng dụng Flutter phục vụ người dùng cuối.

Sau khi lựa chọn phần cứng, khóa luận xây dựng định dạng gói tin LoRa, quy trình gửi nhận dữ liệu, cơ chế xác nhận lệnh và mô hình dữ liệu trên ThingsBoard. Hệ thống được kiểm thử theo từng khối chức năng trước khi kiểm thử tích hợp toàn bộ luồng từ cảm biến đến ứng dụng và từ ứng dụng quay lại thiết bị chấp hành. Kết quả thử nghiệm được phân tích theo ba nhóm tiêu chí: khả năng thu thập và hiển thị dữ liệu theo thời gian thực, khả năng truyền lệnh điều khiển hai chiều, và các hạn chế liên quan đến hiệu chuẩn cảm biến, độ ổn định truyền thông cũng như khả năng mở rộng.
